\documentclass[11pt]{article}

\usepackage[margin=1.15in]{geometry}
\usepackage{amsmath,amssymb,amsthm}
\usepackage{graphicx}
\usepackage{booktabs}
\usepackage{xcolor}
\usepackage[round]{natbib}
\usepackage[colorlinks=true,linkcolor=blue!60!black,citecolor=blue!60!black,urlcolor=blue!60!black]{hyperref}
\usepackage{microtype}

% ---------- theorem environments ----------
\newtheorem{theorem}{Theorem}
\newtheorem{proposition}{Proposition}
\newtheorem{lemma}{Lemma}[section]
\newtheorem{corollary}{Corollary}[section]
\newtheorem{conjecture}{Conjecture}
\newtheorem{observation}{Observation}
\newtheorem*{restatedproposition}{Restatement}
\newtheorem*{restatedconjecture}{Restatement}
\newtheorem*{restatedobservation}{Restatement}
\theoremstyle{definition}
\newtheorem{assumption}{Assumption}
\theoremstyle{remark}
\newtheorem{remark}{Remark}[section]

% Manual-tag assumption blocks used inside appendix proofs (drafts carry
% their own per-proposition assumption labels A1..An).
\newenvironment{assumptionlist}{\begin{description}}{\end{description}}

\newcommand{\pim}{\pi_{\mathrm{m}}}
\newcommand{\pir}{\pi_{\mathrm{r}}}
\newcommand{\popp}{p_{\mathrm{opp}}}
\newcommand{\ktype}{k_{\mathrm{type}}}
\newcommand{\kpiv}{k_{\mathrm{piv}}}

\title{Clearing Conditional Commitments:\\
\large Privacy, Specification, and Custody in a Clearinghouse for Interdependent Action\thanks{%
Preliminary working draft (v1, June 2026). Comments welcome. The numerical companion --- a
chain of nine verified ``toy model'' memos, the solver scripts, and the adversarial
verification logs --- is available from the author. \emph{AI-use disclosure:} generative
AI (Anthropic's Claude) assisted in drafting and typesetting this manuscript from the
author's research notes and proposition drafts, in formalizing and cross-checking the
proofs, and in adversarial verification of the results; the author developed the model and
results, reviewed the entire manuscript, and is responsible for all content. All errors are
mine.}}
\author{Timothy J. Miano\thanks{Affiliate, Massachusetts Institute of Technology; and Founder \& CEO, Proffer Futures, Inc. Contact: \texttt{timmiano@mit.edu}. \emph{Disclosure:} the author holds equity in Proffer Futures, Inc., which operates the ifwishlist platform discussed in Section~9.}}
\date{June 2026 --- preliminary and incomplete}

\begin{document}
\maketitle

\begin{abstract}
\noindent
Latent willingness of the form ``I will, if enough of the right others will'' has no
standard clearing institution. We model a \emph{conditional-commitment clearinghouse}: a platform
that authenticates participants (accuracy $q$), verifies that their mutual conditions are
jointly satisfiable (consensus integrity $c$), routes compatible counterparts to one another,
and chooses three design dials --- custody of principal $\gamma$, specification timing
$\sigma$, and disclosure granularity $k$ --- before activating the coalition at a fixed
point. Our central result is that the welfare-maximizing clearinghouse is \emph{minimal on
all three dials at once}: it holds no principal, fixes terms late, and discloses progress
only through a $k$-anonymized aggregate. The result is stated with its boundary: it holds in
a bounded high-authenticity, intermediate-integrity region, reverses on each axis outside
it, and the three minimalisms are coupled in welfare magnitude rather than additively
separable. Three further results discipline the disclosure claim. First, the privacy dial's
screening cost is a genuine information-leakage object --- the Shannon mutual information
between coalition type and the $k$-anonymized support count --- and full disclosure does not
perfectly screen. Second, the masking advantage survives a refinement-free selection: under
Carlsson--van~Damme risk dominance no rationality parameter appears anywhere in the
model, the revealed game's contested-regime collapse is exact, and masked momentum is
sustained --- proven at the deadline rung and for the smallest coalitions, conjecture-gated
(and grid-verified) at interior rungs. Third, partial
masking is \emph{constrained-optimal for a structural reason}: a single disclosure dial
necessarily bundles ``this coalition is real'' with ``you are not pivotal,'' so --- in a
one-round theorem extended by a dynamic calibration --- the interior optimum exists even
where the unconstrained ideal would be full type-disclosure with zero pivotality leakage. We state three formal propositions with complete proofs, enumerate the
open conjectures, and propose an administrable safe-harbor rule keyed to the mechanism's
structure.\\[4pt]
\textbf{JEL:} C72, D82, D47, H41.\quad
\textbf{Keywords:} conditional commitment, coordination, information design, $k$-anonymity,
market design, assurance contracts.
\end{abstract}

% =====================================================================
\section{Introduction}\label{sec:intro}

\subsection{The vocabulary, fixed first}

The broad problem family is \emph{interdependent action}: situations where each
participant's willingness to act depends on who else will act, in what number, under what
terms, and with what visibility. Collective action, collective consumption, assurance,
focal-point coordination, and matching and scheduling are subcases --- the family is broad,
and we will not pretend any one branch is the genus.

The economic object this paper introduces is the \emph{conditional commitment}: ``I will,
if a qualifying composition of others will.'' Not a pledge (a scalar added to a total), but
a node with edges --- a private commitment whose condition refers to the \emph{composition}
of the committing set: rosters, roles, types, compound thresholds.

The institution that makes the object clearable is the \emph{conditional-commitment
clearinghouse}: private conditional commitments are collected, checked for joint
satisfiability, and activated at the moment the edge-constraints are simultaneously
satisfied --- a fixed point whose output is common knowledge of a feasible coalition. It
clears commitments, not prices. Markets clear goods; this institution clears conditional
commitments --- and in doing so makes previously unclearable interdependent action
actionable.

\subsection{The friction: why the clearinghouse never existed}

Conditional willingness has always existed. No institution cleared it because four barriers
stood at once, and removing any three does not open the institution:
credibility (``I would'' is cheap talk; signals are forgeable), compatibility (mutual
conditions may not be jointly satisfiable --- is there a real coalition, or sixty unrelated
``ifs''?), discoverability (compatible counterparts cannot find each other), and legality
(turning activation into value movement can trigger custody, transmission, or liability
regimes). The institution's non-existence was overdetermined: each wall alone is fatal. The
thesis is that the four walls are now jointly removable for the first time --- credibility
and discoverability by algorithmic inference at scale, compatibility by a mechanism-design
move, legality by a legal-architecture move.

\subsection{The contribution}

Three results, in increasing order of surprise.

\emph{First, a clearable primitive.} We formalize the conditional commitment as a graph
object --- nodes, compositional edge-constraints, a fixed-point activation event --- and
show that a pledge-threshold mechanism is its scalar degenerate case
(Section~\ref{sec:model}).

\emph{Second, one central result, three minimalisms.} In a single welfare objective over custody
$\gamma$, specification $\sigma$, and disclosure $k$, the joint optimum is the minimal
corner on all three axes together --- no custody, fuzzy-early specification, partial
masking --- in a bounded high-$q$, intermediate-$c$ region, reversing off it on each axis
(Section~\ref{sec:central}). The efficient clearinghouse clears \emph{because} it
withholds.

\emph{Third, the disclosure dial decomposed.} The privacy dial controls two channels with
opposite signs --- type/screening (favoring disclosure) and pivotality/protection (favoring
masking) --- and because one physical dial bundles them, partial masking is
constrained-optimal even where the unconstrained ideal is full type-disclosure
(Sections~\ref{sec:dial}--\ref{sec:selection}). The screening cost of masking is quantified
by an actual information measure, and the protection advantage survives a refinement-free
equilibrium selection (with the momentum bound conjecture-gated at interior rungs, as the
proposition statement discloses). Proposition~\ref{prop:band} establishes the interior masking band in
integrity; Proposition~\ref{prop:selection} the refinement-free selection;
Proposition~\ref{prop:bundle} the bundled-dial structure of partial masking.

The displacement claim in one line: prior work either fixes the good, conditions on a sum,
or treats common knowledge as exogenous; we endogenize the good, the coalition membership,
and the disclosure boundary, and then clear relational conditional commitments.

% =====================================================================
\section{Related literature}\label{sec:lit}

Six lineages bracket the contribution; none contains it.

\paragraph{Collective action and provision points.}
\citet{olson1965logic} names the free-rider problem but offers no clearing institution.
\citet{bagnoli1989provision} and the assurance-contract mechanisms of
\citet{tabarrok1998private} fix the good ex ante and condition on a \emph{scalar}
threshold; we endogenize the good's terms and make the condition \emph{relational}
(compositional edge-constraints). \citet{marxmatthews2000dynamic} is the direct dynamic
ancestor --- an observable running total as the Markov coordinating state in voluntary
contribution --- and supplies the completion structure, but contains no disclosure-design
result. \citet{halac2020raising} reach unique implementation through contingent
\emph{payments} to heterogeneous investors over a scalar threshold: the transfer route to
the destination our disclosure route reaches.

\paragraph{Coordination and common knowledge.}
Schelling's focal points and \citet{chwe2001rational} explain when coordination succeeds
given common knowledge; we \emph{manufacture} common knowledge as a service, with
endogenous membership. The global-games selection machinery we draw on is
\citet{carlsson1993global} and \citet{frankelmorrispauzner2003equilibrium}. We cite the
latter's \S6.4 for the criterion's pedigree, not as a blanket license: where the stage game
has strategic complementarities, their \S6.4/Theorem~4 identifies the Laplacian
(uniform-belief) action as the noise-independent selection; at the pivotal (diagonal)
states that drive our result the stage game is a volunteer's dilemma --- commit and wait
are strategic \emph{substitutes} there --- which falls outside the licensed class, so at
those states we adopt the Laplacian action as the \emph{definition} of stagewise risk
dominance, cross-validated by an independent $\lambda$-free potential-maximizer refinement
(Proposition~\ref{prop:selection} and its Appendix~\ref{app:P2} scoping).

\paragraph{Information design.}
\citet{morris2002social} and \citet{cornand2008publicity} establish that public disclosure
precision can flip welfare; \citet{li2023global} studies obfuscation directly. In all
three, precision is a Gaussian noise parameter on a fundamental --- not a quantified
privacy-mechanism leakage measure, and there is no coalition screening rate.
\citet{kamenica2011bayesian}, \citet{bergemann2019information},
\citet{morris2024implementation}, and \citet{inostroza2025adversarial} supply the
persuasion, BCE, and adversarial-selection toolkit over payoff states: lineage, not
preemption. \citet{vosooghi2017information} is nearest by setting --- information design
inside coalition formation --- but the lever is beliefs about a scalar social-cost state,
not a privacy mechanism; its published dynamic extension
\citep{vosooghi2022information} has a third-party sender choosing threshold revelation of
a scalar fundamental to a single coalition, with no privacy mechanism over participant
structure and no screening rate. \citet{halac2021rank} is the closest
privacy-as-instrument neighbor (rank uncertainty defeating strategic uncertainty in team
production), with no quantified leakage primitive. \citet{chen2023stability} is the
matching-side precedent for stability being information-structure-relative, with pairwise
types and no disclosure welfare sign-flip.

\paragraph{Crowdfunding and matching.} The crowdfunding design literature
\citep{strausz2017crowdfunding,ellman2019optimal,belleflamme2014crowdfunding} studies
all-or-nothing thresholds and moral hazard over a \emph{fixed} good with a scalar target;
matching with contracts \citep{hatfield2005matching} clears relational structure but over
fixed contract terms with no disclosure design. Both condition on what we endogenize.

\paragraph{Market design.}
The program of \citet{roth1999redesign,roth2002engineer} --- designing clearing
institutions for market-like environments that did not previously clear --- is the lineage
we claim membership in. \citet{weyl2018radical} is nearest in ambition, distinct in
object: we clear non-transferable conditional commitments, not assets or votes.

\paragraph{The new ingredient,} absent from each lineage by construction, is a
\emph{quantified privacy-leakage primitive} (the mutual information between coalition type
and a $k$-anonymized support count) \emph{driving a coalition screening rate}, joined to a
difficulty-keyed masking advantage. That conjunction is the contribution.

% =====================================================================
\section{The environment}\label{sec:model}

$N$ agents; one latent outcome whose terms are not yet fixed. Each agent $i$ holds a
private type and a conditional commitment $c_i = $ ``act if composition $S_i$ of others
commits.'' A resolver authenticates nodes with accuracy $q$. A router surfaces the outcome
to compatible agents. Coalitions are of two types, good (deliverable) and bad (illusory),
with prior $\pi_b$ on bad; consensus integrity $c$ indexes the probability that the posted
conditional structure is jointly satisfiable.

The platform chooses three dials:
\begin{itemize}
\item \textbf{Custody} $\gamma \in [0,1]$ --- the fraction of principal held at commitment
time, carrying a trust-liability cost $c_{\mathrm{cust}}\gamma$ and deadweight $d_\gamma$,
against a possible commitment-device benefit $\beta_\gamma$ that can lower the effective
activation cutoff.
\item \textbf{Specification} $\sigma \in [0,1]$ --- how early the outcome's terms harden.
Precise-early kills latent commitments via the compatibility factor
$\mathrm{comp}(\sigma)$; fuzzy-too-long adds indeterminacy cost $\mathrm{indet}(\sigma)$.
The outcome's terms may be an \emph{endogenous output} co-resolved by the committed set:
scheduling-by-coalition is the clean instance, where the date is selected by the coalition
rather than fixed before the coalition exists.
\item \textbf{Disclosure} $k \in [0,1]$ --- the granularity of progress disclosure,
implemented as a $k$-anonymity dial on the support count (additive blur
$r(k) = \mathrm{round}((1-k)N)$).
\end{itemize}

Activation occurs when the edge-constraints are simultaneously satisfiable --- the fixed
point --- after which campaign terms harden and value, where present, moves on the
recipient's own rail (the platform holds nothing). The welfare objective of the unified
numerical model is
\begin{equation}\label{eq:W}
W(\gamma,\sigma,k;\,q,c) \;=\;
\mathrm{comp}(\sigma)\,\bigl(1-\rho_{\mathrm{eff}}(k,\gamma,c)\bigr)\,
\bigl(V-\mathrm{indet}(\sigma)\bigr)\,(1-d_\gamma\gamma)
\;-\;\bigl(1-\mathrm{comp}(\sigma)\bigr)\bigl(1-s(k)\bigr)L
\;-\;c_{\mathrm{cust}}\gamma,
\end{equation}
where $\rho_{\mathrm{eff}}$ is the effective unraveling risk and $s(k)$ the bad-coalition
screening rate. We solve for when an activatable coalition forms --- the fixed point exists
\emph{and is selected} --- and for the dial setting that maximizes welfare: coalitions that
should form, do; spurious ones do not.

\paragraph{Reduced-form status.} The unified objective \eqref{eq:W}, the revealed game's
pivotality blend, and the bundled-dial channel coupling are disciplined reduced-form
primitives, not objects derived from a Bayesian extensive form; the appendices state this
where each is used, and each appendix's local model is self-contained (local symbols are
local). The conclusions most robust to this level of abstraction are the structural ones
--- the two-channel decomposition, the bundling constraint, the regime splits; the
quantitative surfaces are measurement targets (Section~\ref{sec:empirics}).

% =====================================================================
\section{The central result: minimal is optimal}\label{sec:central}

\begin{theorem}[informal; numerically established]\label{thm:minimal}
In a bounded high-authenticity, intermediate-integrity region, and under a bounded
custody-benefit, the welfare-maximizing clearinghouse reduces custody, specification, and
disclosure together: the joint optimum $(\gamma^*,\sigma^*,k^*)$ of \eqref{eq:W} is
$\gamma^*=0$ (no custody), $\sigma^*=0$ (fuzzy-early), and $k^*<1$ (partial masking). The
optimum reverses off the region on each axis. The three dials are coupled in welfare
magnitude --- quasi-separable in threshold structure, not additively separable --- with
specification$\times$disclosure the leading interaction.
\end{theorem}

Each apparent weakness does real work. The clearinghouse clears precisely because it
withholds: holding principal adds a trust liability the clearing does not require; forcing
terms early kills latent commitments before the coalition can form; revealing granular
pivotal state can unravel the conditional consensus it is meant to advertise.

Theorem~\ref{thm:minimal} is established numerically (a restricted static disclosure facet
is the part proven analytically, as Proposition~\ref{prop:band} below); four qualifiers travel with
it. \emph{Region, not everywhere}: the joint-all-minimal bands are $c\in[0.66,0.72]$ at
$q=0.95$ and $c\in[0.64,0.74]$ at $q=1.0$, and the region requires $q \gtrsim 0.95$; the
masking band in $c$ is an interior band --- transparency wins below it (screening
dominates) and, for $q<1$, above it (at $q=1$ with the reference parameters the upper
closure fails and masking persists toward $c=1$, Appendix~\ref{app:P1}). \emph{Coupled, not separable}:
the three-dimensional joint argmax equals the product of one-dimensional argmaxes in only
221 of 600 parameter draws; the welfare cost of the separable policy is typically small
(median $0.43\%$) with a real tail ($p_{95}=92\%$). \emph{Custody-minimalism is the
bounded claim}: $\gamma^*=0$ holds only when the commitment-device benefit is small
($\beta_\gamma \ge 0.25$ flips it), the qualifier \citet{diamond1984delegated} already
suggests. \emph{Minimal disclosure means partial masking}: the optimal dial is an interior
$k$-anonymity level, not a blackout; and the \emph{dynamic} (live progress-bar) version of
the masking advantage is a finite-$N$ effect that vanishes in the continuum --- the durable
large-population result is the static coordination/screening logic. The dynamic effect is
real for small campaigns, which is the regime the live platform actually runs.

% =====================================================================
\section{The disclosure dial decomposed}\label{sec:dial}

The single dial $k$ controls two information channels with opposite signs.

\subsection{The type/screening channel}

Disclosure carries a signal about coalition \emph{type} --- good (deliverable) versus bad
(illusory). The channel's precision is grounded in an actual information measure: the
support count $m \sim \mathrm{Binom}(N, p_\psi)$ with $p_G > p_B$ is published under
$k$-anonymity, and precision is the Shannon mutual information $I(\text{type};
m_{\mathrm{obs}})$. The numerical findings, each surviving cross-mechanism
re-implementation: $I(k)$ is monotone (data processing) and \emph{saturating}
($I(1)=0.62$ bits $< H=0.97$ at the reference calibration --- the exact count does not
perfectly reveal type, because the binomial supports overlap); therefore full disclosure
does not perfectly screen, $s(1)<1$, \emph{derived} from the support gap rather than
posited, with $s(1)\to1$ recovered in the large-gap perfect-separation limit; and the
screening rate $s(k)$ is non-monotone under a pessimistic prior. One scope restriction: at
$k=1$ every privacy mechanism collapses to the exact count, so the endpoint is pinned by
$(p_G,p_B)$ alone, and the curve's interior shape is mechanism-dependent. The grounding
closes the endpoint, not the curve; $(p_G,p_B)$ is itself the next-deeper primitive.

\subsection{The pivotality/protection channel}

Disclosure also reveals \emph{pivotal structure} --- who is and is not needed. The two
channels live in complementary excludability regimes: the type channel has bite only at
$\xi < 1-\kappa/V$, where the activation payoff exceeds the free-ride payoff; the
protection channel operates in the complementary regime where free-riding pays
($\xi > 1-\kappa/V$). There, in a contested market, revealing pivotality lets
inframarginal movers free-ride and collapses clearing; masking forces every committer to
act as if pivotal, sustaining the coalition.

\subsection{The static band: a theorem}

The static interplay of screening and unraveling produces the central disclosure
comparative static: partial masking is optimal on a nonempty set contained strictly inside
the interior of the integrity range --- a single band in every tested calibration, with
connectedness the one open conjecture. Figure~\ref{fig:band} displays it; Proposition~\ref{prop:band} proves it.

\begin{figure}[t]
\centering
\includegraphics[width=0.62\textwidth]{figures/fig-band.pdf}
\caption{The interior masking band. Optimal disclosure $k^*(c)$ in the static model at
$q=0.95$ ($V=1$, $L=1.2$, $\tau=0.7$, Beta integrity family with concentration $\eta=8$).
Full disclosure is optimal at both integrity extremes; partial masking on the interior
band (dashed edges). The lower transition is a discontinuous global-max switch; the upper
transition is continuous (Appendix~\ref{app:P1}).}
\label{fig:band}
\end{figure}

\input{appendix/P1-statement}

The proof (Appendix~\ref{app:P1}) characterizes the band sharply for the Beta family:
the first-order--condition violation set is an exact interval, via strict log-concavity of
the marginal-balance ratio (a trigamma inequality); the upper threshold is the root of an
explicit equation with $k^*(c)\to1$ continuously; and Topkis comparative statics give that
raising authentication accuracy $q$ expands the masking set and lowers every optimal
disclosure level. Two scope restrictions: at $q=1$ with the reference parameters the upper
closure fails --- masking persists toward $c=1$, consistent with the numerics --- so the
$c_{\mathrm{hi}}<1$ half is scoped to $q<1$; and connectedness of the masking set is
Conjecture~\ref{conj:connected} (the lower transition occurs while the corner first-order
condition still holds, so no FOC characterizes it).

% =====================================================================
\section{Selection and the structure of partial masking}\label{sec:selection}

Two results complete the disclosure story: the masking advantage needs no equilibrium
refinement parameter, and partial masking has two structural sources.

\subsection{A refinement-free selection}

Earlier numerical versions selected the dynamic equilibrium by a logit rule with
rationality parameter $\lambda$ --- and the masking advantage inverted at high $\lambda$,
leaving the result refinement-dependent. The resolution replaces logit with
Carlsson--van~Damme risk dominance: a $\popp$-consistent Laplacian belief, uniform over
opponents' commit \emph{propensity}, routed through the binomial movement kernel. At
diagonal (volunteer's-dilemma) states, where complementarities fail, the Laplacian action
is adopted as the \emph{definition} of stagewise risk dominance rather than invoked as an
FMP theorem (\S\ref{sec:lit}). No rationality parameter appears anywhere in the model, so
no inversion is possible.

\begin{figure}[t]
\centering
\includegraphics[width=0.62\textwidth]{figures/fig-adv.pdf}
\caption{Masked versus revealed clearing probability across market difficulty $\theta$,
at the reference calibration. The revealed game collapses exactly ($\pir=0$) in the
contested regime; masked momentum is sustained throughout. The masking advantage at
$\theta\le0.30$ is $+0.99$.}
\label{fig:adv}
\end{figure}

\input{appendix/P2-statement}

The contested-regime collapse --- part (iv) --- is exact and fully proven: when
free-riding pays ($\xi > 1-\kappa/V$) and the market is contested, the pivotality weight
classifies every mover as inframarginal, revelation triggers the free-ride action, and
$\pir = 0$ identically. The momentum floor is proven conditional on diagonal ignition
(property DI), which is itself proven outright at the deadline rung and unconditionally
for $K=2$; at interior rungs for $K\ge3$ it is Conjecture~\ref{conj:DI}, verified
pointwise at every calibration grid point. The precise reading of the advantage: strictly
positive when contested, near-tie elsewhere (the numerics give $-8\times10^{-4}$ at
$\theta=0.40$ --- ``$\ge 0$ everywhere'' is false and we do not claim it).

\subsection{The bundled dial}

Why is the optimal masking \emph{partial}? With both channels live in the dynamic game on
the single bundled dial, a genuine interior optimum exists (Figure~\ref{fig:Gk}).

\begin{figure}[t]
\centering
\includegraphics[width=0.62\textwidth]{figures/fig-Gk.pdf}
\caption{The bundled-dial inverted-U\@. Activation gain $G(k)$ at the reference contested
calibration ($\theta=0.30$, $\xi=0.667$, $T=10$): both endpoints are dead --- full masking
cannot ignite momentum, full disclosure fully unravels --- and the interior peak near
$k=0.70$ carries the entire welfare range.}
\label{fig:Gk}
\end{figure}

\input{appendix/P3-statement}

The decoupled-dial result --- part (iii) --- is what makes the interior structural rather
than incidental: with two separate dials, the optimum is the corner
$(\ktype,\kpiv)=(1,0)$, full type-disclosure with zero pivotality leakage. The bundled
interior is therefore a \emph{constrained} optimum created by the yoking --- one physical
dial cannot certify a coalition's reality without leaking its pivotal structure --- and
never a type-versus-pivotality balance. Partial masking thus has two distinct sources:
the integrity-distribution shape (Proposition~\ref{prop:band}) and the one-dial coupling
(Proposition~\ref{prop:bundle}).

Horizon scope: parts (ii)--(iii) are $T=1$ theorems; the endpoint results hold at
every horizon. At $T\ge2$ the original general-horizon decoupled-corner conjecture was
\emph{falsified} in adversarial verification --- within the proposition's own assumptions,
the corner can fail to be attained, and positive pivotality leakage can strictly improve
welfare at fixed type-disclosure (Observation~\ref{obs:leakagehelps}): a leakier future
cuts the waiting option and re-ignites mid-states through the continuation channel.
Whether the decoupled argmax always retains $\kpiv=0$ is open
(Conjecture~\ref{conj:decoupledT}). The leakage-helps phenomenon is itself a new,
empirically probeable prediction.

% =====================================================================
\section{Boundary: where the theorem reverses}\label{sec:boundary}

The theorem is conditional on the walls actually being down.

\emph{$q$ collapses.} Authentication is weak; masked disclosure now shields
\emph{manufactured} commitments rather than protecting a real coalition. ``Minimal
disclosure'' flips from efficient to fraud-enabling. Below the located threshold
(the joint region requires $q\gtrsim0.95$), transparency dominates.

\emph{$c$ collapses.} The conditional structure is illusory --- nodes are real but their
conditions are not jointly satisfiable; clearing produces no activatable outcome. The
masking band is interior in $c$ (Proposition~\ref{prop:band}): below it, screening
dominates and transparency wins; above it, coalitions are too robust to unravel and
disclosure is harmless.

\emph{The contested regime is absent.} The protection channel matters only where the
static selection frontier binds and free-riding pays ($\xi > 1-\kappa/V$): a good
coalition that exists but would unravel if its pivotal structure were exposed. Where
completion is easy, masking is neutral (the near-tie region of
Proposition~\ref{prop:selection}).

\emph{Endogeneity over-extends.} If outcome terms stay fuzzy past the binding moment, the
campaign instrument is indeterminate; precision must harden as commitment strength rises
(the interior $\sigma^*$ under high indeterminacy cost).

\emph{The taxonomy over-expands.} ``Interdependent action'' becomes a slogan unless every
example reduces to a conditional-commitment graph --- nodes, edges, constraints, fixed
point. The family is broad; the primitive must stay narrow.

% =====================================================================
\section{Mechanism and an administrable rule}\label{sec:rule}

\paragraph{Mechanism.} A no-custody clearinghouse for conditional commitments with
(a)~authenticated nodes, (b)~verifiable conditional structure, (c)~endogenous outcome
terms with a hardening schedule, (d)~partial disclosure (an interior $k$-anonymity level)
on granular state, (e)~value riding the recipient's own rail, and (f)~non-transferable
participant commitments.

\paragraph{Normative rule.} We propose model-motivated design criteria for a limited
safe harbor for a privacy-preserving conditional-commitment clearinghouse --- a proposed
administrable rule, not a legal conclusion, and pending counsel review. The four criteria:
non-transferable participant commitments; no custody of principal; an auditable
node-authentication process; and small-$N$ disclosure controls. The rule regulates the
\emph{clearing of commitments} --- not raw wishes, raw data, platform funds, or
transferable financial instruments --- and is keyed to the conditional-commitment
structure, not the campaign's vertical. The theory motivates each prong: the
authentication prong is the $q$-collapse boundary; the disclosure-controls prong is the
small-$N$ regime where the dynamic protection effect (and its abuse surface) lives; the
no-custody and non-transferability prongs keep the cleared object a coordination device
rather than a financial instrument.

% =====================================================================
\section{Implementation: a constructive implementability case}\label{sec:constructive}

The empirical spine is not yet a dataset --- it is a live platform's release history,
readable as a constructive case that the institution is implementable. Implementability is
all this section claims: a release history cannot validate optimality, welfare, or the
selection results, and the calibration quantities are unmeasured
(Section~\ref{sec:empirics}). The
platform (ifwishlist, live June 2026) did not ship unrelated features; it removed the four
walls one at a time: authenticated counterpart resolution (credibility, $q$); a
conditional-commitment grammar with compound and roster thresholds (compatibility, $c$);
semantic routing (discoverability); a no-custody rail (legality); and the $k$-anonymity
dial with a specification-hardening schedule --- an architecture consistent with the
model's recommended design.

% =====================================================================
\section{Empirical agenda}\label{sec:empirics}

The theory chain has largely converged; the marginal value has moved to measurement. Four
quantities, all currently unmeasured, each first-order for a specific result:

\begin{center}
\begin{tabular}{lll}
\toprule
\textbf{Quantity} & \textbf{Governs} & \textbf{Disciplines} \\
\midrule
support gap $p_G - p_B$ & the screening cap $s(1)$, large-gap reversal & \S\ref{sec:dial} \\
excludability $\xi$ & which channel regime governs & \S\ref{sec:dial}--\ref{sec:selection} \\
campaign-size distribution & the scale of the finite-$N$ protection effect & \S\ref{sec:central} \\
$\pi_b$ and the integrity distribution & band location; interior $k^*$ & Prop.~\ref{prop:band}, \ref{prop:bundle} \\
\bottomrule
\end{tabular}
\end{center}

A further theory iteration is warranted only when a measured quantity makes a specific
model question first-order. Observation~\ref{obs:leakagehelps} adds a sharp new item:
the leakage-helps regime is directly probeable on live activation data.

% =====================================================================
\section{Scope, open questions, and the verification record}\label{sec:scope}

Every numerical claim in this paper comes from a chain of solver scripts under fixed
seeds, and every analytical result was drafted against those scripts as ground truth and
then attacked by three separately instantiated adversarial verification passes --- a
step-by-step logic audit, a counterexample hunt running randomized parameter sweeps inside
the stated assumption sets, and a scope audit of claim--proposition alignment --- each run
with no access to the others' findings, with scripts and logs archived alongside the
numerical companion. Roughly
$3{,}500$ randomized counterexample instances were run against the stated assumption sets;
no proved statement was refuted. The process corrected the first draft of every artifact
it touched --- including four sections of this paper --- and falsified one of our own
conjectures, producing Observation~\ref{obs:leakagehelps}. We report the open residue explicitly:

\begin{enumerate}
\item \textbf{Conjecture~\ref{conj:connected}} (connectedness of the masking band): the
lower transition is a discontinuous global-max switch with no first-order
characterization.
\item \textbf{Conjecture~\ref{conj:DI}} (diagonal ignition at interior rungs): closing it
needs a simultaneous induction bounding the on-cone waiting value.
\item \textbf{Conjecture~\ref{conj:interiorT}} (the $T>1$ interior): numerically robust;
the band grows with horizon.
\item \textbf{Conjecture~\ref{conj:decoupledT}} (decoupled argmax retains zero leakage):
open after refinement; the pointwise version is false
(Observation~\ref{obs:leakagehelps}).
\item Deriving $(p_G,p_B)$ from a deeper compatibility primitive; correlated signals; an
endogenous adverse-selection prior $\pi_b$.
\end{enumerate}

Two further scope notes. The legality-wall claim inherits open predicates pending counsel
review (the no-custody rail's money-transmission posture; charitable-solicitation
regimes). And the live system does not yet instrument $q$ or $c$; both are required for
the central theorem to bite on data.

% =====================================================================
\section{Conclusion}\label{sec:conclusion}

AI does not make the engine intelligent. It removes two of the four walls --- credibility
and discoverability --- that kept interdependent willingness from ever clearing. The
mechanism removes the third (compatibility); the legal architecture removes the fourth.
With all four down, ``I would, if\ldots'' becomes a clearable economic object --- and the
efficient clearinghouse is the one that holds nothing, fixes nothing early, and hides the
granular.

\bibliographystyle{plainnat}
\bibliography{references}

\clearpage
\appendix

\section{Proof of Proposition~\ref{prop:band}}\label{app:P1}
\input{appendix/P1-proof}

\clearpage
\section{Proof of Proposition~\ref{prop:selection}}\label{app:P2}
\input{appendix/P2-proof}

\clearpage
\section{Proof of Proposition~\ref{prop:bundle}}\label{app:P3}
\input{appendix/P3-proof}

\end{document}
